\documentclass[11pt,a4paper]{article}
\usepackage[margin=2.5cm]{geometry}
\usepackage{amsmath,amssymb,amsthm}
\usepackage{booktabs}
\usepackage{graphicx}
\usepackage[hidelinks]{hyperref}
\usepackage{xcolor}

\newcommand{\R}{\mathbb{R}}
\newcommand{\bu}{\mathbf{u}}
\newcommand{\bv}{\mathbf{v}}
\newcommand{\bF}{\mathbf{F}}
\newcommand{\bn}{\mathbf{n}}
\newcommand{\rhat}{\hat{\mathbf{r}}}
\newcommand{\Pc}{\mathcal{P}}
\newcommand{\Pt}{\widetilde{\mathcal{P}}}
\newcommand{\Kc}{\mathcal{K}}
\newcommand{\code}[1]{\texttt{#1}}
\newtheorem{proposition}{Proposition}

\title{Staggered kite finite elements for the shallow-water equations\\[2pt]
\large Notes on the implementation in \code{dmplex.jl}}
\author{Oswald Knoth}
\date{Working notes, \today}

\begin{document}
\maketitle

\begin{abstract}
These notes document the staggered kite discretization of the rotating
shallow-water equations (EGU 2026 poster) as implemented in the Julia package
\code{dmplex.jl}: the kite grid and its cell centers, the finite-element spaces
and their degrees of freedom, the discrete operators, the conservation
properties of the linear and nonlinear schemes, and numerical results. Open
questions --- the lowest-order pairing, other velocity elements, the pair of
Wess et al.\ with Coriolis, damping and parallelization --- are collected at
the end.
\end{abstract}

\tableofcontents

%------------------------------------------------------------------------------
\section{Equations}

On a closed surface $\Omega$ (the unit sphere or a doubly periodic plane) with
outward unit normal $\rhat$, the rotating shallow-water equations in
vector-invariant form are
\begin{align}
  \partial_t \bu + q\,\rhat\times\bF + \nabla\bigl(K + g h\bigr) &= 0, &
  \partial_t h + \nabla\cdot\bF &= 0,
\end{align}
with mass flux $\bF = h\bu$, kinetic energy $K = \tfrac12|\bu|^2$, absolute
vorticity $\omega = \zeta + f$, $\zeta = \rhat\cdot\nabla\times\bu$, and
potential vorticity $q = \omega/h$ (so $q\,\rhat\times\bF = \omega\,\rhat\times\bu$).
The equations conserve mass $\int h$ and energy
$E = \int \tfrac12 h|\bu|^2 + \tfrac12 g h^2$. For the sphere we use units with
radius $a = 1$ and rotation rate $\Omega = 1$, so $f = 2z$.

%------------------------------------------------------------------------------
\section{The kite grid}

\paragraph{Construction.} Let $\Pc$ be a polygonal primal mesh (triangles,
quadrilaterals, the hexagons and pentagons of a Voronoi mesh). Every cell
$P\in\Pc$ with vertices $w_1,\dots,w_n$ (counter-clockwise seen from outside)
is split into $n$ quadrilaterals (\emph{kites})
\[
  K_i = (w_i,\; m_i,\; c_P,\; m_{i-1}),\qquad i = 1,\dots,n,
\]
where $c_P$ is the cell center and $m_i$ a point on the primal edge
$(w_i,w_{i+1})$. The kites sharing a primal vertex $w$ form the dual cell
$\widetilde P_w$. In \code{create\_kite\_mesh} every kite has the canonical
local numbering: vertex~1 is the primal vertex, vertex~3 the center; local
edges~1 and~4 are halves of primal edges (\emph{half edges}), local edges~2
and~3 join an edge point to the center (\emph{interior edges}); the interior
edges of $P$ are exactly the boundaries of the dual cells. Labels record the
parent cell, the dual cell and the vertex and edge types.

\paragraph{Cell centers.} The center is chosen so that it coincides with the
vertex of the orthogonal dual mesh whenever there is one
(\code{cell\_centers}):
\begin{itemize}
\item triangles: the circumcenter (vertex of the Voronoi dual);
\item Voronoi meshes: the generators. Every Voronoi edge with great-circle
  normal $\bn_e = a\times b$ is the perpendicular bisector of the generators
  $g_i, g_j$ of its two cells,
  \begin{equation}
    g_i - g_j \parallel \bn_e, \qquad \bn_e\cdot(g_i + g_j) = 0,
  \end{equation}
  which is a homogeneous linear system for all generators. Its least-squares
  null vector is computed by sparse inverse iteration. For MPAS meshes this
  recovers \code{xCell} to $10^{-14}$, for the ICON dual \code{vlon/vlat}
  to $10^{-14}$;
\item other meshes (cubed sphere, HEALPix, planes): the centroid.
\end{itemize}
On meshes with an orthogonal dual, $m_i$ is the crossing of the primal edge with
the dual edge $(c_P, c_{P'})$, so the interior kite edges lie on the dual edges
and each dual cell is exactly the dual polygon (the Delaunay triangle of an MPAS
mesh); otherwise $m_i$ is the edge midpoint. The number of kites equals
$\sum_P n_P$, e.g.\ $245\,760$ for the MPAS mesh \code{x1.40962}.

%------------------------------------------------------------------------------
\section{Finite-element spaces}

All kite spaces live on the reference square $[-1,1]^2$ with the primal vertex
at $(-1,-1)$ and the center at $(1,1)$. The reference map is the bilinear
interpolant of the four kite vertices, radially projected onto the sphere
(curved elements); a chordal variant without projection is available.

\paragraph{Height: \code{KiteRadau(p)}.} $Q_p$ on every kite with nodes at the
tensor Gauss--Radau points $r_0<\dots<r_p = 1$ that include the end point at the
center side. Nodes lie in the kite interior ($p^2$), on the interior edges
($p$ each) and at the center (one), never on primal edges. Hence $h$ is
continuous inside a primal cell and broken across primal edges,
\[
  V_h = \{ h \in L^2(\Omega) : h|_P \in H^1(P)\ \forall P\in\Pc,\ h|_K\in Q_p \},
\]
with $1 + n_P(p + p^2)$ degrees of freedom per primal cell. With the Radau nodes
as quadrature the mass matrix is diagonal; the rule is exact for degree $2p$
per direction, so for $p\ge1$ the lumped mass integrates the bilinear
area element exactly and is consistent. For $p = 0$ the single Radau node sits
at the center corner and $4\,\det J(c)$ is an $O(1)$ wrong kite area (kites do
not become parallelograms under refinement); since $h$ is then constant per
primal cell the exact mass matrix (the cell area) is used instead.

\paragraph{Velocity: \code{KiteBrokenRT(k)}.} $\mathrm{RTCF}_k$ on every kite
with normal continuity across half edges (inside a dual cell) and none across
interior edges (the dual cell boundaries):
\[
  V_u = \{ \bu \in L^2(\Omega)^2 : \bu|_{\widetilde P} \in H(\mathrm{div},\widetilde P)
  \ \forall \widetilde P\in\Pt,\ \bu|_K\in\mathrm{RTCF}_k \}.
\]
Edge degrees of freedom are normal components at the $k+1$ Gauss points; the
degrees of freedom of the two interior edges are stored with the kite itself.
The exact mass matrix is block diagonal with one block per dual cell (9 for
RT0 on a triangular dual cell, 12 for a quadrilateral one, 30 for RT1 on a
triangular dual cell).

\paragraph{Other velocity spaces.} Two alternatives with the same continuity
pattern (normal components continuous across half edges, broken across
interior edges) are implemented:
\begin{itemize}
\item \code{KiteBrokenBDM(k)}: $\mathrm{BDM}_k$ on every kite,
  $[P_k]^2\oplus\mathrm{span}\{\mathrm{curl}(x^{k+1}y),\mathrm{curl}(xy^{k+1})\}$,
  normal components at the $k+1$ Gauss points of each edge and, for $k\ge2$,
  moments against $[P_{k-2}]^2$; 8 instead of 12 degrees of freedom per kite for
  $k=1$. The interior-edge degrees of freedom are stored with the kite as for RT.
\item \code{KiteDualRadau(p)}, the velocity space of the dual cell method of Wess,
  Kapidani, Codecasa and Sch\"oberl (2024): $Q_p\times Q_p$ on every kite with nodes at
  the tensor Gauss--Radau points that include the primal-vertex side (the mirror
  image of the height nodes). On the half edges the degrees of freedom are the
  normal components, shared by the two kites of the dual cell; inside the kite the
  component values. There are no nodes on the interior edges, so the field is
  broken there without any special storage. With the nodes as quadrature the mass
  matrix couples only the components at one node: blocks of size 3 for $p=1$ on
  hexagon kites, independent of the mesh.
\end{itemize}

\paragraph{Vorticity.} Continuous $Q_{k+1}$ (\code{Lagrange(k+1)}) on the kites
with Gauss--Lobatto nodes and lumped, diagonal mass.

\paragraph{Implementation.} A \code{Section} assigns degrees of freedom to mesh
points (vertices, edges, cells). Layouts that depend on point labels extend
\code{num\_nodes(layout, dm, p)}, and an element may store dofs on another point
than the one they belong to geometrically (the interior-edge RT dofs on the
kite). The closure of a kite then delivers the global indices, with node order
reversed and signs flipped on reversed edges, so the spaces are conforming
where intended; the tests check the continuity pattern across both edge types
to machine precision.

%------------------------------------------------------------------------------
\section{Discrete operators}

With the contravariant Piola map
$\bv = J\hat\bv/\det J$ one has $\bv\cdot\nabla\eta\,dA = \hat\bv\cdot\hat\nabla\hat\eta\,d\xi$,
$(\nabla\cdot\bv)\,\eta\,dA = (\hat\nabla\cdot\hat\bv)\,\hat\eta\,d\xi$ and
$\bv\cdot\bn\,ds = \hat\bv\cdot\hat\nu\,dt$ with the scaled reference edge normal
$\hat\nu$; both operators below are therefore independent of the metric.

\paragraph{Divergence.} For $\eta\in V_h$, $\bF \in V_u$,
\begin{equation}
  D(\eta;\bF) = \sum_{K} \int_K \bF\cdot\nabla\eta
  - \sum_K \sum_{e\subset\partial K\cap\partial\Pc} \int_e \eta_K\,\bF\cdot\bn_K,
\end{equation}
summed over the half edges of each kite with the trace of $\eta$ from inside
$K$; as a matrix $B_{ij} = D(\eta_i;\bv_j)$.

\paragraph{Gradient.} For $p\in V_h$, $\bv\in V_u$,
\begin{equation}
  G(p;\bv) = \sum_K \int_K p\,\nabla\cdot\bv
  - \sum_K \sum_{e\subset\partial K,\ e\ \text{interior}} \int_e p\,\bv_K\cdot\bn_K,
\end{equation}
where $p$ is single valued on interior edges.

\begin{proposition}
$G(p;\bv) = -D(p;\bv)$, i.e.\ $G = -B^T$, and $\mathbf 1^T B = 0$.
\end{proposition}
\begin{proof}
Integrating by parts on each kite,
$\int_K p\,\nabla\cdot\bv = -\int_K \nabla p\cdot\bv + \int_{\partial K} p\,\bv\cdot\bn_K$.
The boundary integral splits into half edges and interior edges; the interior
edges cancel against the face term of $G$, leaving
$-\sum_K\int_K\nabla p\cdot\bv + \sum_K\sum_{\text{half}}\int p_K\,\bv\cdot\bn_K = -D(p;\bv)$.
For $\eta \equiv 1\in V_h$, $D(1;\bv) = -\sum_K\sum_{\text{half}}\int\bv\cdot\bn_K = 0$,
since $\bv\cdot\bn$ is continuous across half edges and both sides appear with
opposite normals.
\end{proof}
Both matrices are assembled independently in the code; the test suite checks
$\|G + B^T\|_\infty < 10^{-12}$.

\paragraph{Vorticity, kinetic energy, Coriolis.}
$\langle\zeta,\gamma\rangle = -\int \bu\cdot(\rhat\times\nabla\gamma)$ for
$\gamma\in Q_{k+1}$ (lumped mass),
$\langle K,\eta\rangle = \int\tfrac12|\bu|^2\eta$ (lumped mass of $V_h$), and
$C_{ij} = \int f\,\bv_i\cdot(\rhat\times\bv_j)$, which is antisymmetric. The
surface normal $\rhat$ is computed from the Jacobian,
$\rhat = J_1\times J_2/\det J$, which is $x/|x|$ on curved elements and the
element normal on chordal ones.

%------------------------------------------------------------------------------
\section{Semi-discrete schemes and conservation}

\paragraph{Linear.}
\begin{equation}
  M_h\dot h = H\,B\,\bu, \qquad M_u\dot\bu = -C\bu - g\,B^T h.
\end{equation}
Mass $\mathbf 1^T M_h h$ is conserved because $\mathbf 1^T B = 0$, energy
$\tfrac12 H\bu^T M_u\bu + \tfrac12 g h^T M_h h$ because $C$ is antisymmetric
and the gradient is the negative adjoint of the divergence.

\paragraph{Nonlinear, poster form.}
\begin{align}
  \langle\dot h,\eta\rangle &= \sum_K\int_K h\bu\cdot\nabla\eta
     - \sum_K\sum_{\text{half}}\int \eta_K\,h_K\,\bu\cdot\bn_K,\\
  \int\dot\bu\cdot\bv &= -\int(\zeta+f)\,(\rhat\times\bu)\cdot\bv + G(gh+K;\bv),
\end{align}
with $h_K$ the trace from inside the kite. On primal edges $h$ is two valued, so
the flux $h_K\,\bu\cdot\bn$ is not single valued: neither mass nor energy is
conserved exactly.

\paragraph{Nonlinear, conservative variant (default).} Project the mass flux,
$\bF = \Pi_{V_u}(h\bu)$, i.e.\ $M_u\bF = \bigl(\int h\bu\cdot\bv_i\bigr)_i$ (one solve
with the block-diagonal velocity mass), and use the potential vorticity
$q = (\zeta+f)/h$ at the quadrature points:
\begin{align}
  M_h\dot h &= B\bF, &
  \int\dot\bu\cdot\bv &= -\int q\,(\rhat\times\bF)\cdot\bv - B^T(gh+K).
\end{align}
\begin{proposition}
The conservative variant conserves mass and the discrete energy
$E_h = \int\tfrac12 h|\bu|^2 + \tfrac12 g\,h^T M_h h$ exactly.
\end{proposition}
\begin{proof}
Mass: $\mathbf 1^T M_h\dot h = \mathbf 1^T B\bF = 0$. Energy: with the same
quadrature in $E_h$ and in the right-hand sides,
$\partial E_h/\partial\bu = M_u\bF$ and
$\partial E_h/\partial h = M_h(K + g h)$. Hence
$\dot E_h = \bF^T M_u\dot\bu + (K+gh)^T M_h\dot h
 = -\int q\,(\rhat\times\bF)\cdot\bF - \bF^T B^T(gh+K) + (gh+K)^T B\bF = 0$,
since $(\rhat\times\bF)\cdot\bF = 0$ pointwise.
\end{proof}
This is the framework of McRae and Cotter (2014) transferred to the kite
spaces; all mass matrices stay diagonal or block diagonal. The test suite
checks $\dot E_h$ and $\mathbf 1^T M_h\dot h$ to round-off on hexagonal,
triangular and quadrilateral kites and on periodic planes.

%------------------------------------------------------------------------------
\section{Numerical results}

Williamson test case~2 (steady geostrophic flow, $u_0 = 0.2$, $g = 1$,
$T = 1$) measures $\|h(T)-h(0)\|_{M_h}$; kites of the hexagonal (Voronoi)
meshes \code{hex(r)} have 960, 3840 and 15360 kites for $r = 2,3,4$.

\subsection{Lowest order}
$L^2$ errors of the discrete gradient $M_u^{-1}B^T\Pi p$ and divergence
$M_h^{-1}B\,\Pi\bu$ against the projections of the exact operators for smooth
$p$ and $\bu$ (hexagon kites with 960, 3840, 15360 kites; cubed-sphere kites with
384, 1536, 6144 kites):

\begin{center}
\begin{tabular}{llccc}
\toprule
pairing & mesh & gradient error & divergence error \\
\midrule
$Q_0\times\mathrm{RT}_0$ & hex  & 7.4, 7.6, 7.7 & 0.017, 0.0069, 0.0030 \\
                         & cube & 5.6, 5.8, 5.9 & 0.065, 0.076, 0.067 \\
$Q_1\times\mathrm{RT}_0$ & hex  & 0.10, 0.052, 0.026 & 0.23, 0.21, 0.21 \\
                         & cube & 0.19, 0.094, 0.047 & 0.40, 0.36, 0.31 \\
$Q_1\times\mathrm{RT}_1$ & hex  & 0.21, 0.11, 0.052 & 0.016, 0.0060, 0.0027 \\
                         & cube & 0.31, 0.15, 0.077 & 0.044, 0.022, 0.012 \\
$Q_2\times\mathrm{RT}_1$ & hex  & 0.0038, 0.00091, 0.00023 & 0.088, 0.045, 0.023 \\
                         & cube & 0.0091, 0.0023, 0.00057 & 0.059, 0.035, 0.019 \\
\bottomrule
\end{tabular}
\end{center}
At lowest order one of the two operators is inconsistent. With $Q_0$ heights the
pressure is constant per primal cell and only sets the half-edge normal
components of the gradient; the kite-local interior-edge velocity degrees of
freedom are not controlled by the pressure. With $Q_1\times\mathrm{RT}_0$ the
discrete divergence of the broken RT0 field tested with $Q_1$ does not converge.
From $\mathrm{RT}_1$ on both operators converge. This agrees with the
experience with CGDycore that the lowest order does not work.

\subsection{Linear model}
$Q_1\times\mathrm{RT}_1$, RK4, geostrophic steady state:

\begin{center}
\begin{tabular}{lccc}
\toprule
primal mesh & coarse & medium & fine \\
\midrule
hexagons (162, 642, 2562 cells) & $3.09\cdot10^{-3}$ & $7.72\cdot10^{-4}$ & $1.94\cdot10^{-4}$ \\
triangles (80, 320, 1280)       & $9.67\cdot10^{-3}$ & $2.46\cdot10^{-3}$ & $6.16\cdot10^{-4}$ \\
cubed sphere (96, 384, 1536)    & $7.88\cdot10^{-3}$ & $1.94\cdot10^{-3}$ & $4.88\cdot10^{-4}$ \\
\bottomrule
\end{tabular}
\end{center}
Second order on all meshes; mass to $10^{-15}$, energy up to the RK4 error.

\subsection{Nonlinear model}
Williamson test case~2, RT1 velocity, CG2 vorticity, SSP-RK3:

\begin{center}
\begin{tabular}{lcccc}
\toprule
variant & $h$ & hex(2) & hex(3) & mass / energy \\
\midrule
poster form          & $Q_1$ & $2.05\cdot10^{-3}$ & $8.76\cdot10^{-4}$ & $\dot M\approx10^{-4}$, $\Delta E/E\approx10^{-4}$ \\
poster form          & $Q_2$ & $5.9\cdot10^{-4}$  & $5.2\cdot10^{-3}$ (grows) & not conserved \\
projected flux       & $Q_2$ & $3.60\cdot10^{-4}$ & $8.79\cdot10^{-5}$ & exact / $10^{-12}$ \\
projected flux + $q$ & $Q_2$ & $3.60\cdot10^{-4}$ & $8.80\cdot10^{-5}$ & exact / exact \\
\bottomrule
\end{tabular}
\end{center}
The pair RT1 / $Q_2$ / CG2 is the most accurate and, with the projected flux,
second-order convergent. With the poster's own-trace flux, $Q_2$ heights grew
on the finer mesh; nothing damps the jumps of $h$ across primal edges there.

\paragraph{Quadrature.} All nonlinear terms use one Gauss rule per kite with
precomputed metric terms ($J/\det J$, $dA$, $\rhat$, $f$). Errors on hex(2) /
hex(3) and cost per SSP-RK3 step on hex(3):
\begin{center}
\begin{tabular}{lccc}
\toprule
rule & hex(2) & hex(3) & ms/step \\
\midrule
$3\times3$ & $3.606\cdot10^{-4}$ & $8.765\cdot10^{-5}$ & 42 \\
$4\times4$ & $3.601\cdot10^{-4}$ & $8.796\cdot10^{-5}$ & 53 \\
$5\times5$ & $3.601\cdot10^{-4}$ & $8.796\cdot10^{-5}$ & 68 \\
\bottomrule
\end{tabular}
\end{center}
The minimal rule ($3\times3$ for RT1) is the default; energy conservation is an
algebraic property and holds for every rule. A time step does not allocate.

\paragraph{Curved versus chordal elements.} With precomputed metric terms both
cost the same; chordal (flat bilinear) kites give $9.10\cdot10^{-4}$ and
$2.50\cdot10^{-4}$ (rate 1.87) instead of $3.60\cdot10^{-4}$ and
$8.80\cdot10^{-5}$ (rate 2.03). Curved elements are the default.

\subsection{Velocity spaces: RT, BDM and the pair of Wess et al.}
Operator consistency ($L^2$ errors of the discrete gradient and divergence on
hexagon kites with 960, 3840, 15360 kites) and the linear geostrophic steady state
with Coriolis ($\|h(1)-h(0)\|$, RK4, finest mesh and rate):
\begin{center}
\begin{tabular}{lccc}
\toprule
pair (velocity / height) & gradient & divergence & geostrophic, finest (rate) \\
\midrule
RT1 / $Q_1$            & 0.21, 0.11, 0.052     & 0.016, 0.0060, 0.0027 & $1.94\cdot10^{-4}$ (1.99) \\
RT1 / $Q_2$            & 0.0038, 0.00091, 0.00023 & 0.088, 0.045, 0.023 & $4.12\cdot10^{-5}$ (2.10) \\
BDM1 / $Q_1$           & 0.13, 0.065, 0.032    & 0.28, 0.26, 0.25      & $1.46\cdot10^{-4}$ (1.78) \\
BDM1 / $Q_2$           & 0.0021, 0.00048, 0.00012 & 0.84, 0.83, 0.82   & $1.61\cdot10^{-4}$ (1.20) \\
Wess $p=1$ ($Q_1^2$ / $Q_1$) & 0.14, 0.069, 0.035 & 0.21, 0.18, 0.18    & $1.68\cdot10^{-4}$ (1.86) \\
Wess $p=2$ ($Q_2^2$ / $Q_2$) & 0.0070, 0.0018, 0.00044 & 0.026, 0.013, 0.0064 & $5.61\cdot10^{-6}$ (2.07) \\
\bottomrule
\end{tabular}
\end{center}
The divergence error of BDM1 and of the Wess pair with $p=1$ does not converge in
this norm (the discrete divergence of the broken velocity tested with the height
space), yet the geostrophic steady state converges for all pairs. The Wess pair
with $p=2$ is the most accurate, about seven times better than RT1/$Q_2$; energy is
conserved exactly by all pairs.

\paragraph{Wess pair with Coriolis.} Assembling the Coriolis term with the velocity
nodes as quadrature (lumped) instead of exactly makes no difference
($1.68\cdot10^{-4}$ versus $1.63\cdot10^{-4}$ for $p=1$, $5.61\cdot10^{-6}$ versus
$5.62\cdot10^{-6}$ for $p=2$ on the finest mesh, same rates). In the nonlinear
Williamson test case~2 (conservative variant, SSP-RK3) the Wess pair with $p=2$
gives $8.54\cdot10^{-5}$ and $2.02\cdot10^{-5}$ on hex(2), hex(3) (rate 2.08)
against $3.61\cdot10^{-4}$, $8.77\cdot10^{-5}$ for RT1/$Q_2$, with mass and energy
conserved; with $p=1$ it converges with rate 1.4. In this implementation the pair
therefore works with Coriolis and in the nonlinear case. Differences to a failing
implementation could lie in the face terms (here the kite divergence $B$ and
$G=-B^T$, both checked independently), in the orientation of the shared normal
components at the half edges, or in the nonlinear terms (here the conservative
variant with projected flux).

\subsection{Kite method versus the standard mixed setup on the kite quads}
The kite grid is an ordinary quadrilateral mesh, so the compatible pair
RT1 $\times$ DG1 $\times$ CG2 of McRae and Cotter (global, non-diagonal mass
matrices, \code{NonlinearShallowWater}) runs on it unchanged. Williamson test
case~2, SSP-RK3, same time step:
\begin{center}
\begin{tabular}{llccccc}
\toprule
mesh & method & $u$ dofs & $h$ dofs & error & mass solve & ms/step \\
\midrule
hex(2) & kite RT1/$Q_2$/CG2   &  9600 &  5922 & $3.61\cdot10^{-4}$ & blocks (30) & 6.9 \\
       & mixed RT1/DG1/CG2    &  7680 &  3840 & $1.94\cdot10^{-4}$ & Cholesky    & 20.0 \\
hex(3) & kite RT1/$Q_2$/CG2   & 38400 & 23682 & $8.77\cdot10^{-5}$ & blocks (30) & 38.4 \\
       & mixed RT1/DG1/CG2    & 30720 & 15360 & $4.72\cdot10^{-5}$ & Cholesky    & 98.5 \\
\bottomrule
\end{tabular}
\end{center}
Both are second order and conserve energy; the mixed setup is about $1.9\times$
more accurate with fewer degrees of freedom. The timings are not a fair
comparison of the methods: the mixed model in \code{dmplex.jl} still uses the
older allocating kernels, and its global mass solves (sparse Cholesky) are the
part that the kite method replaces by local block inverses, which matters most
for explicit time stepping, on GPUs and in parallel.


\subsection{Galewsky jet}
2562 hexagons, 15360 kites ($\approx 182$\,km), $\Delta t = 66.9$\,s, 6 days:
the jet becomes unstable after day~3 and rolls up by day~6; mass changes by
$6\cdot10^{-12}$, energy by $1.2\cdot10^{-8}$ (time integration only). Grid-scale
noise is visible near the vortices (no dissipation).

\subsection{Bickley jet}
Doubly periodic $[0,4\pi)^2$, $u = \mathrm{sech}^2(y-2\pi)$, no rotation, Froude
number 0.2, perturbation in wavenumber~1; $32\times32$ quadrilaterals, 4096 kites,
$t = 40$: vortex street with two cat's-eye pairs; mass $8\cdot10^{-11}$, energy
$8\cdot10^{-10}$.

%------------------------------------------------------------------------------
\section{Open questions and proposals}

\subsection{Other velocity elements}
Broken BDM and the space of Wess et al.\ are implemented (Section~3 and the
results). Open: BDM2 and $p\ge3$ for the Wess pair, the dual-cell method on
triangles only (Wess et al.\ use triangular primal meshes) versus general
polygons, and a lumped version of the RT velocity mass.

\subsection{Damping}
For comparison, Gusto's standard shallow-water model (\code{SIQNModel}) obtains
its dissipation implicitly from upwinded DG transport (\code{DGUpwind}) of the
depth and of the velocity in vector-invariant form, with SSP-RK3 transport inside
a semi-implicit quasi-Newton step; interior-penalty diffusion and SUPG are
optional. Options for the kite method, in order of preference:
\begin{enumerate}
\item \emph{APVM} (anticipated potential vorticity): replace $q$ by
  $q - \tau\,\bu\cdot\nabla q$, $\tau = \alpha\Delta t$, in the vorticity term.
  The term stays of the form $q^*\,\rhat\times\bF$, so energy is still conserved
  exactly while potential enstrophy is dissipated (McRae and Cotter 2014).
  $\nabla q$ is available from $\nabla\zeta$ (CG2) and $\nabla h$ inside the kite.
\item \emph{Upwind height flux on primal edges}, where $h$ jumps: replace the
  own-trace $h_K$ by the upwind trace in the poster form. The flux is single
  valued, so mass is conserved, and energy is dissipated by the jumps of $h$
  (the analogue of Gusto's upwinded depth transport).
\item \emph{Biharmonic hyperviscosity} on the velocity,
  $\Delta\bu = \nabla\delta - \rhat\times\nabla\zeta$ with the discrete
  divergence $\delta = M_h^{-1}B\bu$ and vorticity $\zeta$, applied twice,
  $\nu_4\propto\Delta x^4/\tau_4$; as in MPAS and CGDycore, mainly for comparison.
\end{enumerate}

\subsection{Comparison with the standard mixed setup on the kite quads}
The kite grid is an ordinary quadrilateral mesh, so the standard compatible pair
RT$_k\times$DG$_k\times$CG$_{k+1}$ (\code{NonlinearShallowWater} on the kite
\code{Plex}) can be run on it directly; its mass matrices are not block
diagonal. See the results section.

\subsection{Parallelization}
\paragraph{Kernels.} The nonlinear terms are KernelAbstractions kernels over a
list of kites: gather the degrees of freedom, evaluate at the quadrature points
with the stored metric terms, and scatter with atomic adds (Atomix). On the CPU
backend the kernels run on the Julia threads, on a GPU unchanged; the scatter
order and therefore the round-off depend on the scheduling. The block-diagonal
mass solves need no atomics since the blocks are disjoint. (The alternative to
atomics is colouring: kites of one colour share no degrees of freedom, colours
are processed one after the other.)

\paragraph{Distribution (MPI).} The primal cells are partitioned with the zonal
equal-area regions of Leopardi (Atlas \code{EqualRegionsPartitioner}, CGDycore
\code{DecomposeEqualArea}); kites follow their primal cell. Each rank holds one
vertex ring of ghost primal cells. Ownership: height dofs follow their primal
cell, velocity dofs their dual-cell block (the smallest part around the dual
vertex), vorticity dofs the smallest part among the kites touching them. With
the ghost ring every owned dof has all its kites locally, so a rank computes
the right-hand sides over its owned and ghost kites, and they are complete for
its owned dofs; all mass blocks of owned dofs are local. Kites and degrees of freedom
are numbered locally (owned first, then ghosts); the global numbering is only used
during setup and to gather states. Per tendency evaluation the halo exchange
updates the velocity $\bu$ and height $h$ (prognostic), the vorticity $\zeta$ and,
in the conservative variant, the mass flux $\bF$ (diagnostic). The kinetic energy
is not exchanged: it lives in the height space, whose ghost degrees of freedom
belong to ghost primal cells whose kites are all local. The linear operators are
applied to the owned rows only. The runs on one and on three ranks agree to
$10^{-12}$, and mass and energy are conserved across ranks. The global mesh is
still built on every rank during setup.

\paragraph{First scaling test.} Galewsky setup, 2562 hexagons (15360 kites),
Mac with 8 performance and 8 efficiency cores, warmed-up SSP-RK3 steps; per
tendency the maximum over ranks:
\begin{center}
\begin{tabular}{rrrrrr}
\toprule
ranks & ms/step & kernels & halo (4$\times$) & sparse products & block solves \\
\midrule
1 & 241 & 34.7 & 0.9 & 5.5 & 19.8 \\
2 & 145 & 21.0 & 1.6 & 4.4 & 10.5 \\
4 &  90 & 14.6 & 1.7 & 3.1 &  6.2 \\
8 &  75 & 13.0 & 3.6 & 2.8 &  4.9 \\
\bottomrule
\end{tabular}
\end{center}
Up to four ranks the speed-up is reasonable (2.7); from four to eight ranks the
kernels hardly get faster, although each rank has half the kites. The
per-kite cost grows by a factor of 2.4, which points to the global-length
vectors (poor locality of the gathers) and to ranks on efficiency cores. With
the local numbering 1, 2 and 4 ranks took 205, 129 and 82\,ms per step; a repeated
test on a loaded machine (load average 9--11) was not meaningful, so the
scaling beyond four ranks still has to be measured on an idle machine.

\end{document}
